% Introduction — Securing LLMs Against Prompt Injection
% Academic manuscript section (IEEE / arXiv style)
% Citation keys match references.txt

\section{Introduction}
\label{sec:introduction}

Large language models are increasingly embedded in chat assistants, developer tools, enterprise copilots, customer-support systems, and agentic workflows that accept free-form natural language as a primary control interface.
This flexibility is also a security liability.
Because models are trained to follow instructions expressed in text, an adversary can often embed conflicting directives in a user prompt and induce the system to ignore application policy, reveal restricted information, adopt an unauthorized persona, execute unintended tool calls, or otherwise pursue an attacker-chosen goal~\cite{owasp_llm,perez2022ignore,yi2024jailbreak}.
Unlike classical injection vulnerabilities that target a brittle parser or query language, prompt injection exploits the model's intended strength---open-ended instruction following---which makes both attack and defense intrinsically linguistic~\cite{11359704,bair2025prompt}.
OWASP identifies this failure mode as prompt injection and ranks it among the most critical risks for LLM applications~\cite{owasp_llm}.
Subsequent work has shown that the threat is not limited to overt phrases such as ``ignore previous instructions'': paraphrased, role-play, hypothetical, multilingual, and optimization-based attacks can achieve similar ends~\cite{yi2024jailbreak,perez2022ignore}, while indirect injection can deliver adversarial content through retrieved documents, web pages, or tool outputs rather than through the primary chat box~\cite{greshake2023indirect,hines2024spotlighting,debenedetti2024agentdojo}.
As surveys emphasize, the underlying difficulty is a collapse of trust boundaries between trusted instructions and untrusted natural-language inputs, turning context integrity into a first-class security problem~\cite{11359704,correia2026systematic,bair2025prompt}.

In practice, many deployments respond first with inexpensive lexical controls: regular expressions, keyword deny-lists, role-token constraints, and prompt-hardening guidance~\cite{openrouter2026promptinjection,getstream2024prompt,mend2024prompt}.
These defenses are attractive because they are transparent, auditable, fast, and easy to update when a new attack signature appears.
Yet the same literature that recommends them also documents their brittleness.
Attackers can preserve malicious intent while changing surface form; multi-step prompts can dilute keyword density; and benign technical language can collide with attack vocabulary, producing false positives that degrade usability~\cite{yi2024jailbreak,wang2025promptsleuth,bair2025prompt,kholkar2025capture}.
The result is a persistent robustness--usability trade-off: tightening rules improves recall on known attacks but harms ordinary users, while loosening rules restores usability at the cost of residual risk.

Stronger alternatives exist across several research clusters.
Structured query interfaces separate instructions from data so that untrusted text is less likely to be treated as executable guidance~\cite{struq2024}.
Preference-optimization methods such as SecAlign internalize resistance by training models to prefer responses to intended instructions over responses to injections~\cite{secalign2024}.
External sanitization and guardrail models detect, rewrite, or remove injected content before backend generation~\cite{shi2025promptarmor,geng2025pisanitizer,chowdhury2025prepsilonepsilonmptsanitizingsensitiveprompts}.
Semantic detectors further argue that unauthorized task intent can remain invariant under paraphrase even when lexical cues disappear~\cite{wang2025promptsleuth}, and context-aware evaluations show that both missed attacks and over-defense must be measured together~\cite{kholkar2025capture,zizzo2025adversarial}.
These lines of work are important, but they do not fully answer a question that arises in many real deployments: when the backend model is fixed, the chat interface cannot be redesigned, and only a portable pre-inference gate can be inserted, how effective is a model-agnostic input sanitizer that combines lexical, structural, and semantic evidence?

This paper addresses that question through a controlled study of context-aware input sanitization for direct prompt injection.
We place a pre-inference gate in front of an instruction-tuned open model and compare five conditions under a shared evaluation protocol: an unprotected baseline, a regular-expression filter, a keyword heuristic, a hybrid handcrafted context-aware sanitizer (Method~A), and a learned soft-fusion variant of the same signal family (Method~B).
Method~A is not framed as a pure ensemble scorer.
It uses a high-confidence hard-trigger layer---especially a semantic veto based on MiniLM similarity to known injections---as the primary defense, while a weighted combination of lexical, structural, semantic, and DistilGPT-2 perplexity cues provides secondary coverage for borderline cases~\cite{reimers2019sentence,wang2020minilm,radford2019language}.
Method~B retains an expanded semantic representation, including maximum similarity, top-three mean similarity, and a benign-contrast term, and replaces handcrafted aggregation with logistic regression and a validation-selected threshold, without hard vetoes.
This contrast lets us ask not only whether context-aware sanitization outperforms lexical baselines, but also whether observed gains come from soft score fusion or from hard semantic gating---and ablations later show that Method~A behaves as a hybrid system in which the latter dominates.

Evaluation is deliberately two-layered, following contemporary guardrail-benchmarking practice that warns against equating filter misses with end-to-end harm~\cite{zizzo2025adversarial,promptfoo2024owasp}.
\emph{Bypass Rate} measures whether the sanitizer fails to block an injection at Layer~1.
\emph{True Attack Success Rate (True ASR)} measures whether an unblocked injection also elicits judged compliance from the target model at Layer~2.
An unblocked injection that the model refuses therefore contributes to Bypass Rate but not to True ASR.
We additionally report false positives on ordinary and edge-case benign prompts, paraphrase and adaptive rewriting robustness, component ablations, and sanitizer latency.
The corpus is compact and curated by design.
The study should therefore be read as a reproducible proof-of-concept comparison under a fixed direct-injection threat model, not as a claim of production completeness against agentic, multimodal, or strongly optimized attack leaderboards~\cite{debenedetti2024agentdojo,yeo2025multimodal,yi2024jailbreak}.

Empirically, lexical defenses reduce Bypass Rate only modestly on the held-out test partition and leave True ASR unchanged relative to the unprotected baseline, indicating that blocking a few known surface forms is not enough once model compliance is measured.
Both context-aware methods improve security substantially.
Method~A achieves the strongest protection, including zero successful attacks on the static test set, but incurs a higher false-positive burden, especially on technical edge-case prompts that reuse attack-like vocabulary in harmless contexts.
Method~B offers a milder usability cost while leaving residual end-to-end risk.
Ablations show that Method~A should be interpreted as a hybrid defense: semantic hard triggers provide most of the protection, while soft weighted scoring is only complementary; Method~B's gains depend primarily on semantic features, and removing all semantic cues returns it to unprotected baseline behaviour on this partition.
Adaptive stylistic rewriting degrades all active defenses but preserves the same ordering, reinforcing that context-aware sanitization remains useful under controlled restatement even though it is not invulnerable.
Latency measurements further show that the added security of Method~A comes with a measurable but moderate CPU-side pre-inference cost concentrated in embedding and perplexity scoring.

The contributions of this paper are as follows:
\begin{enumerate}
  \item We formulate direct prompt-injection defense as a model-agnostic pre-inference sanitization problem and implement Method~A as a hybrid gate in which semantic hard triggers form the primary layer and weighted multi-signal scoring provides secondary coverage.
  \item We compare this hybrid handcrafted design (Method~A) against learned soft fusion without hard vetoes (Method~B), enabling an evidence-based attribution of security gains to architectural choices rather than to pipeline complexity alone.
  \item We evaluate all methods under a two-layer protocol that separates Bypass Rate from True ASR, and we complement the main comparison with false-positive, paraphrase, adaptive-rewriting, ablation, and latency analyses.
  \item We show empirically that lexical baselines are insufficient on this corpus once end-to-end compliance is measured, that Method~A's defensive capability originates mainly in semantic hard triggers rather than ensemble scoring alone, and that Method~B's gains are driven primarily by semantic features under an explicit usability trade-off.
\end{enumerate}

The remainder of the paper is organized as follows.
Section~\ref{sec:related-work} reviews prompt-injection taxonomies, rule-based and alignment defenses, sanitization and semantic detection, and evaluation practice.
Section~\ref{sec:methodology} presents the threat model, dataset, defense designs, and two-layer evaluation methodology.
Section~\ref{sec:experimental-setup} details the experimental protocol, model configurations, robustness settings, and reproducibility controls.
Section~\ref{sec:results} reports the main results, robustness findings, ablations, and latency measurements.
Section~\ref{sec:limitations} states the scope of the study and its limitations.
Section~\ref{sec:conclusion} concludes.
